% Complete proofs for the 2001 source-native pattern families.
\section{Evocation and implication patterns: proofs}
\label{sec:p18-erotetic-patterns-proofs}

Use the source languages, abbreviations, finite families and syntactic
conditions of Section~\ref{sec:p18-erotetic-patterns-and-composition}.
Evocation requires both coverage and every individual non-entailment;
implication requires coverage from each initial answer and a nonempty
proper original-answer witness for every further answer. The source states
the examples but does not print these full proofs
\cite[secs.~6,8]{Wisniewski2001QuestionsInferences}.

In descriptions of concrete models below, assigning a displayed domain subset
to P or R specifies that predicate's extension in that model.
Within object-language formulas, P and R remain predicate symbols;
the model-description set role is local and does not turn Px into set membership.

\begin{annexitem}{annex:p18.ep.propositional-evocation}{thm:p18-ep-propositional-evocation}{Every propositional example and every unsettled answer}
\begin{proof}
Unmentioned variables can be assigned arbitrarily. For (6.5), every valuation
has a true member of \(\{p,\neg p\}\). The premise is true at both
\(p=0\) and \(p=1\), which respectively falsify its two answers.
For (6.6), polar coverage is again automatic. The premise \(p\lor q\)
is true at \((p,q)=(0,1)\), falsifying \(p\), and at \((1,0)\),
falsifying \(\neg p\). For (6.7), those same two premise-valuations
respectively falsify \(\neg q\) and \(q\). Thus both polar questions
are covered without either of their answers being entailed.

For (6.8), \(p\supset q\) is true at \((0,0)\) and \((1,1)\).
These falsify \(p\) and \(\neg p\); polar coverage holds as before.
For (6.9), the premise is true at \((0,0)\) and \((0,1)\), falsifying
\(q\) and \(\neg q\). For (6.10), a true biconditional has exactly
the truth pairs \((0,0)\), \((1,1)\). In the first, the negative
conjunctive answer is true and the positive one false; in the second,
the positive answer is true and the negative one false. These prove
coverage and both required non-entailments.

For (6.11), truth of \(p\lor q\) gives a true answer \(p\) or \(q\).
The pairs \((0,1)\), \((1,0)\) respectively falsify those answers.
For (6.12), its premise permits exactly \((1,1),(1,0),(0,1)\). Their
true answers are respectively \(p\land q,p\land\neg q,\neg p\land q\),
which proves coverage. For each proposed answer choose either of the
other two pairs; the premise remains true and that answer is false.
This checks all three non-entailments.

For (6.13), the premises \(p\supset(q\lor r),p\) imply \(q\lor r\),
so cover \(\{q,r\}\). The triples \((p,q,r)=(1,0,1)\) and
\((1,1,0)\) satisfy both premises and respectively falsify \(q,r\).
For (6.14), if \(\neg r\) and \((p\land q)\supset r\) are true,
\(p,q\) cannot both be true. Thus \(\neg p\) or \(\neg q\) is true.
The triples \((1,0,0)\) and \((0,1,0)\) satisfy both premises and
respectively falsify \(\neg p\), \(\neg q\).
For (6.15), a true \((p\lor q)\lor r\) means either \(p\) is true
or at least one of \(q,r\) is true, giving the stated two-answer coverage.
The triples \((0,1,0)\), \((1,0,0)\) satisfy the premise and respectively
falsify \(p\), \(q\lor r\).

All displayed answers are syntactically distinct for the stipulated distinct
variables, and every list has at least two members. Each case therefore
meets both evocation clauses and question well-formedness.
\end{proof}
\end{annexitem}

\begin{annexitem}{annex:p18.ep.firstorder-polar}{thm:p18-ep-firstorder-polar}{Models for both responses to every first-order polar question}
\begin{proof}
Each polar question is covered by the ordinary classical truth clauses,
including the truth clauses for closed quantified sentences and equality.
It remains to show that each of its two answers is not entailed.

For (6.16), take the one-element model with every constant naming its
element and \(P\) true of that element. The observations and
\(\forall xPx\) are true, so the negative answer is not entailed.
Take instead a domain with two elements, put every displayed constant at
the first, and let \(P\) hold just of that first element. The observations
remain true but the universal is false, so the positive answer is not entailed.

For (6.17), a one-element model with both predicates everywhere true makes
all observations and the universal conditional true. A two-element model
with the displayed constants at 0, \(P=\{0,1\}\), and \(R=\{0\}\)
makes all observations true but the universal conditional false at 1.
These witness non-entailment of both polar answers.

For (6.18), give \(c_1,\ldots,c_{n+1}\) distinct values in a domain
of \(n+1\) elements. Let \(P\) hold of all those elements and let
\(R\) hold of the first \(n\). Every conjunctive observation and
\(Pc_{n+1}\) is true, but \(Rc_{n+1}\) is false. Extending \(R\)
to include the last element keeps the premises true and makes the target
positive answer true. Thus neither target answer follows. Distinct
denotations here are a permissible countermodel choice, not a standing
assumption on all models.

For (6.19), interpret \(c_1,c_2\) at the same P-element in a one-element
domain. Both observations and the identity are true. Interpret them at
different P-elements in a two-element domain; both observations remain
true but the identity is false. Equality is neither settled by the
observations nor treated as identity of the two constant expressions.
This supplies both required non-entailments in every case.
\end{proof}
\end{annexitem}

\begin{annexitem}{annex:p18.ep.named-evocation}{thm:p18-ep-named-evocation}{All listed and unlisted constant-answer cases}
\begin{proof}
For (6.20), choose a witness to \(\exists xPx\). It either equals the
interpretation of \(c\), in which case \(Pc\) is true, or differs from
it, in which case \(\exists x(Px\land x\ne c)\) is true. This proves
coverage without requiring the witness to have a name. In domain
\(\{0,1\}\) with \(c=0\), the interpretation \(P=\{0\}\)
satisfies the premise and falsifies the second answer; \(P=\{1\}\)
satisfies it and falsifies the first. Thus neither answer is entailed.

For (6.21), a witness to \(\exists x(Px\land N_{\vec c}(x))\)
equals the interpretation of at least one \(c_i\); hence \(Pc_i\)
is true. For any candidate answer \(Pc_i\), interpret the displayed
constant symbols distinctly and let \(P\) hold only at another listed
constant's value, possible since \(n>1\). This satisfies the premise
and falsifies that answer. No unique-names principle is needed for
coverage, while distinct denotations are allowed in the countermodels.

For (6.22), the negative universal has a witness making
\(N_{\vec c}(x)\) true and \(Px\) false. Equality gives some false
\(Pc_i\), hence the required negative-answer coverage. To falsify
the candidate \(\neg Pc_i\), use distinct displayed denotations,
make \(P\) true at \(c_i\)'s value and false at another listed value.
Then the negative universal is true and that negative answer is false.

For (6.23), the coverage proof for (6.21) supplies a true member of the
larger wh-answer set as well. Now let \(c\) be any constant of the
language. Interpret distinctly the finite set of constant symbols
consisting of \(c\) and the displayed list. If \(c\) is one of
the displayed symbols, choose a different listed symbol; if it is not,
choose any listed symbol. Let \(P\) hold only at that chosen listed
value. The premise is true and \(Pc\) false. This argument covers
every possible wh-answer, including every name outside the finite list.

For (6.24), the coverage proof for (6.22) likewise supplies a true
wh-answer. For any candidate \(\neg Pc\), interpret the same finite
set of needed symbols distinctly, make \(P\) true at \(c\)'s value,
and false at a listed value different from it. Such a listed value exists
by \(n>1\) if \(c\) is listed, and trivially if it is not. Assign
other predicate values arbitrarily. The negative universal is true
and the candidate negative answer false. Thus no wh-answer is entailed.

All countermodels have finite nonempty domains. Other constants, if
infinitely many exist, can be interpreted at any one of their elements;
only the finitely many needed distinct symbols were given distinct values.
Every explicit question has its required number of syntactically distinct
closed answers. Both evocation clauses are therefore established in all
five examples at their actual answer-set extents.
\end{proof}
\end{annexitem}

\begin{annexitem}{annex:p18.ep.complete-list}{thm:p18-ep-complete-list}{Coverage and non-entailment for arbitrary complete-list answers}
\begin{proof}
Let a model satisfy the displayed finite disjunction and choose a true
\(H_{\vec c^{j}}\). At each listed denotation, its naming disjunction
is true, so the biconditional makes \(P\) true there. At any P-individual,
the biconditional makes that naming disjunction true. Hence every positive
conjunct and the exhaustion conjunct of \(W_{\vec c^{j}}\) are true.
That is a direct answer to \(?U(Px)\), proving coverage.

Fix any proposed direct answer \(W_{\vec t}\). Take the finite set of
constant symbols occurring in all the displayed alternative lists and
\(\vec t\), use those symbols as distinct elements of a finite domain,
and interpret each at its corresponding element. Interpret other constants
arbitrarily. Let \(C_j\) be the denotation set of list
\(\vec c^{j}\) and \(T\) the denotation set of \(\vec t\).
The alternative sets \(C_j\) are nonempty and pairwise disjoint, so
at least two are different. Consequently \(T\) cannot equal all of them.
Choose \(j\) with \(T\ne C_j\), and give \(P\) extension exactly
\(C_j\). Then \(H_{\vec c^{j}}\) and the whole premise are true.

In this model \(W_{\vec t}\) could be true only if \(T\subseteq P\)
by its positive conjuncts and \(P\subseteq T\) by its exhaustion
clause, that is, only if \(T=C_j\). Since these sets differ, either
some listed positive instance is false or the exhaustion clause is false.
Thus the candidate answer is not true. The argument applies to every
finite candidate list, including permutations, lists mixing alternative
names, single-name lists, and lists using other names. It proves every
individual non-entailment, not merely non-entailment of the displayed
candidate answers.
\end{proof}
\end{annexitem}

\begin{annexitem}{annex:p18.ep.evocation-boundaries}{thm:p18-ep-evocation-boundaries}{Explicit failures under altered standing conditions}
\begin{proof}
If every individual is required to equal one of the observed \(c_i\),
the premises of (6.16) make \(P\) true of every individual, so entail
the positive universal answer. The premises of (6.17) similarly make
\(R\) true of every P-individual and entail the positive universal
conditional. Evocation fails its non-entailment clause in each case.
Adding \(c_1\ne c_2\) to the equality inquiry explicitly entails its
negative answer, so that modified basis does not evoke it.

With just one name, the finite lists in (6.21),(6.22) have one answer
and are not questions. In the wh-versions, a witness to
\(\exists x(Px\land x=c_1)\) entails \(Pc_1\); a witness to
\(\neg\forall x(x=c_1\supset Px)\) entails \(\neg Pc_1\).
Hence (6.23),(6.24) would fail their informativeness clause. A single
\(H_{\vec c}\), without alternative extension disjuncts, entails
\(W_{\vec c}\) by the already supplied positive/exhaustion argument.
Thus that altered complete-list premise does not evoke its question.

Finally, in a two-element domain interpret every constant at 0 and
\(P\) only at 1. The bare existential is true but every named answer
\(Pc\) is false. It does not even cover \(?S(Px)\). The restricted
existential in (6.23) is false in this model, since its naming disjunction
does not hold at the only P-individual. This explains exactly which added
condition supplied the named witness without changing the overall model
class to one with names for every individual.
\end{proof}
\end{annexitem}

\begin{annexitem}{annex:p18.ep.pure-patterns}{thm:p18-ep-pure-patterns}{Every pure direction, truth cell and singleton limit}
\begin{proof}
For (8.1), every initial answer \(\neg A\), \(\neg\neg A\)
covers the polar target. Target answer \(A\) entails original answer
\(\neg\neg A\); target answer \(\neg A\) entails itself.
For (8.2), target \(\neg A\) entails original \(\neg A\), and
target \(\neg\neg A\) entails original \(A\). Polar coverage again
holds. Each original list has two elements, so these singleton witnesses
are nonempty proper. For (8.3), both questions have the same answer set.
An initial answer covers itself in the target; each target answer entails
itself in the original. This proves the stated regular implications.

For (8.7), write the four binary answers as the cells \(11,10,01,00\)
according to the A- and B-values they express. The labels are abbreviations
for those four distinct conjunction formulas, not additional answers.
Every model has exactly one true cell. The target polar question is always
covered. For each target response choose the following cell set:
\[
\begin{array}{c|c|c}
 \otimes& A\otimes B&\neg(A\otimes B)\\\hline
 \land&\{11\}&\{10,01,00\}\\
 \lor&\{11,10,01\}&\{00\}\\
 \supset&\{11,01,00\}&\{10\}\\
 \equiv&\{11,00\}&\{10,01\}.
\end{array}
\]
Each set is nonempty and proper in the four-answer source list. The
corresponding response makes a cell from precisely that set true, by the
connective's truth clause. These are all the required reduction witnesses,
whether or not all four truth pairs are realizable by the chosen formulas.

For (8.10), every initial answer covers the binary question since every
model has a true cell. Each target cell entails either \(A\) or
\(\neg A\), according to its first sign, giving a singleton original
witness. For (8.11), take the singleton determined by the second sign.
This proves regularity in both cases without independent-atom hypotheses.

For the claimed nonregular cases in (8.7), take A and B as independent
atoms. For conjunction use the negative response, whose three truth cells
are realized by three valuations. For disjunction and implication use
their positive responses, again realized in three different cells. For
equivalence use the positive response, realized in cells 11 and 00.
No single binary answer is true in all those valuations: the cells are
mutually exclusive and at least two of them occur. Thus that target
response entails no individual original answer, although the displayed
proper subset is entailed. This refutes regularity for each connective.
\end{proof}
\end{annexitem}

\begin{annexitem}{annex:p18.ep.general-refinement}{thm:p18-ep-general-refinement}{Arbitrary formula refinement and coarsening}
\begin{proof}
For the refinement, original answer A covers itself in the target;
original answer \(B\lor C\) covers one of B,C by the Boolean
disjunction clause. Target answer A entails original A. Target answers
B and C each entail original \(B\lor C\). These singleton witnesses
are proper in the two-answer original list, proving regularity.

For the reverse, original answers A,B,C each cover A or \(B\lor C\).
Target A has singleton witness \(\{A\}\). Target \(B\lor C\)
has witness \(\{B,C\}\), which is nonempty and proper in the
three-answer original list. Thus ordinary implication holds, for
all the formulas meeting the displayed syntax conditions.

For analyticity of the refinement, every immediate child of target A
is already a subformula of original A; every immediate child of B or
C is a subformula of the original answer \(B\lor C\). For
coarsening, every immediate child of A is a subformula of original A,
and the children B,C of the target disjunction are themselves original
direct answers. Each direction is therefore analytic.

To refute regularity of coarsening in general, let A,B,C be distinct
atoms. The valuations \((A,B,C)=(0,1,0),(0,0,1)\) both make
target \(B\lor C\) true. A is false in both; B and C are each
false in one. No individual original answer is entailed by that target.
For a dependent instance where coarsening is regular, let
\(A=r,B=p,C=\neg\neg p\), with distinct atoms p,r. The two
questions remain syntactically well formed. The target disjunction
entails B, while A entails itself, providing singleton witnesses.
Hence the statement is a nonregularity possibility, not a prohibition
on every formula instance. No atomic independence condition was used
to prove the two general valid implications.
\end{proof}
\end{annexitem}

\begin{annexitem}{annex:p18.ep.analytic-profile}{thm:p18-ep-analytic-profile}{Syntactic, not semantic, analyticity}
\begin{proof}
In (8.1), every immediate child of target answer A is a subformula of
the source answer \(\neg A\), and the child A of target \(\neg A\)
is also such a subformula. In (8.2), the children of target \(\neg A\)
and \(\neg\neg A\) are A and \(\neg A\), both original direct
answers or their subformulas. In (8.3), every child of either target
answer is a subformula of that same formula in the original answer set.
Together with their proved pure implications, these prove analyticity.

In (8.7) for conjunction, the immediate children of \(A\land B\)
are A,B, and the child of \(\neg(A\land B)\) is \(A\land B\).
All are subformulas of the first binary answer. For another connective,
the target negative answer has immediate child \(A\otimes B\).
This compound is larger than each of A,B, so cannot be a subformula
inside either. Outside those components, a binary-answer subformula can
only be one of its four outer conjunctions, or \(\neg A,\neg B\).
The compound has neither conjunction nor negation as its outermost sign,
so is none of these. It is also not a negation of any subformula, since
its outermost sign is not negation. Thus it lies outside
\(\mathcal H(?\pm|A,B|)\), and analyticity fails. This reasoning
uses finite syntactic parsing, not independence or Boolean equivalence.

In (8.10), the immediate children of the four target answers are exactly
among \(A,\neg A,B,\neg B\), and all four occur. The first two
are already in \(\mathcal H(?A)\); therefore the analytic condition
is exactly membership of both B and \(\neg B\) there. Interchanging
the two components proves the analogous condition for (8.11).
For independent distinct atoms, the other atom is not an original
subformula or its negation, so the condition fails.

For the first positive counterinstance take \(A=p\land q,B=p\).
Then p is a subformula of A and \(\neg p\) its allowed negation.
Consequently both required memberships hold, so (8.10) is analytic.
For (8.11), interchange these A,B choices. They are syntactically
distinct formulas; every binary answer remains a distinct formula.
These instances show why the generic non-analytic classification cannot
be read as forbidding analyticity in every compound-formula instance.
\end{proof}
\end{annexitem}

\begin{annexitem}{annex:p18.ep.signed-cofactors}{thm:p18-ep-signed-cofactors}{Every fixed sign and both possible target responses}
\begin{proof}
The target is polar, hence covered after any initial answer. It suffices
to give the regular witness for each of its responses. If A is fixed true
or false, respectively, the compound truth function as a function of B is
\[
\begin{array}{c|c|c}
 \otimes&A\text{ true}&A\text{ false}\\\hline
 \land&B&0\\
 \lor&1&B\\
 \supset&B&1\\
 \equiv&B&\neg B
\end{array}
\]
Here 0 and 1 denote constant truth values, not new primitive formulas.
The conditional truth clause gives value 1 whenever its antecedent is false.
For each response B or \(\neg B\), these eight cases specify a
constant truth value of the principal compound. Select the singleton
containing that compound if its value is 1 and its negation if 0.
The fixed sign and the response entail that answer. Each singleton is
proper in the two-answer principal list.

If B is fixed true/false, the corresponding four pairs of functions of A
are conjunction \(A,0\), disjunction \(1,A\), conditional
\(1,\neg A\), and biconditional \(A,\neg A\). For each response
A or \(\neg A\), choose the principal answer by the same 1/0 rule.
This checks each of the eight fixed-B cases as well. These are direct
truth-clause deductions under the fixed premise; a source question-chain
composition is not used.

For example, on \(\{\neg A\}\) the conjunction is false before
either B-response is received. Both responses have the same negative
principal answer as witness. With independent atoms B is still unsettled
by that basis, but the initial compound is not. More generally no
informativeness clause is part of this implication proof.
\end{proof}
\end{annexitem}

\begin{annexitem}{annex:p18.ep.polar-refinements}{thm:p18-ep-polar-refinements}{Every covered-list index and the three-answer case}
\begin{proof}
For (8.14), target polar coverage is automatic. At a valuation where A
is true, either B or \(\neg B\) is true, giving a member of
\(\{A\land B,A\land\neg B\}\). At a valuation where \(\neg A\)
is true, that original answer itself is true. These witnesses are
nonempty and proper in the three-element original answer set. For
independent atoms, target A is compatible with B true and with B false;
neither conjunctive original answer holds in both, and \(\neg A\)
holds in neither. Hence it does not entail an individual original answer.

For (8.15), polar target coverage again holds regardless of which
initial answer is true. Target \(A_i\) entails its singleton original
answer. With target \(\neg A_i\), the coverage premise \(D_n\)
forces some \(A_j\) with \(j\ne i\) true. This gives the stated
complement witness, which is nonempty because \(n>1\), and proper
because it excludes \(A_i\). For \(n=2\) it is a singleton.

If \(n>2\) and the original answers are independent atoms, fix
target \(\neg A_i\). For candidate \(A_i\) it is already false.
For any other candidate \(A_j\), choose \(k\notin\{i,j\}\), make
only \(A_k\) true, and all other original atoms false. The coverage
premise and target response are true but the candidate answer is false.
Thus that response entails no individual original answer. These
countervaluations do not weaken the already proved multiple-answer witness.
\end{proof}
\end{annexitem}

\begin{annexitem}{annex:p18.ep.denial}{thm:p18-ep-denial}{Coverage, every reduction witness and both indispensable premises}
\begin{proof}
A valuation satisfying \(\neg K_n\) makes at least one \(A_i\)
false, so covers the negative target list. In particular this holds
with every initial answer added to the full basis. If the same
valuation satisfies target \(\neg A_i\) and \(D_n\), at least one
different \(A_j\) is true. Thus the stated complementary answer set
is entailed in the multiple-conclusion sense. It is nonempty proper;
for \(n=2\) it is a singleton, proving regularity.

For larger independent-atom lists, fix response \(\neg A_i\).
For any candidate \(A_j\) with \(j\ne i\), choose a third index
k and make just \(A_k\) true. Both basis formulas and the response
are true while that candidate is false. Candidate \(A_i\) is false
by the response itself. Therefore regularity need not hold.

Omitting \(\neg K_n\), the valuation making every original atom
true satisfies the remaining basis and every initial answer, but
no target negative answer is true. The first implication clause fails.
Omitting \(D_n\), the valuation making every original atom false
satisfies \(\neg K_n\) and every target response, but no original
answer is true. It refutes every possible second-clause witness.
These failures identify separately the work of each premise.
\end{proof}
\end{annexitem}

\begin{annexitem}{annex:p18.ep.partition-patterns}{thm:p18-ep-partition-patterns}{Every finite block split, both clauses and regularity limits}
\begin{proof}
Write the two original answer blocks as \(\{A_j:j\leq i\}\) and
\(\{A_j:j>i\}\). They are nonempty and proper because
\(1\leq i<n\). For (8.18), the polar target is always covered.
Its response B and the first conditional force \(D^-_i\), hence
a true answer in the first block. Response \(\neg B\) and the
second conditional force \(D^+_i\), hence a true answer in the second.

For (8.19), the basis formula \(B\lor C\) covers the two-answer
target with any initial answer. Target B and its conditional force a
first-block answer; target C and its conditional force a second-block
answer. For (8.20), any first-block initial answer makes \(D^-_i\)
true and its biconditional gives B. Any second-block initial answer
similarly gives C. Conversely B with the first biconditional forces
the first block, and C with the second forces the second block.
These prove both clauses in all three patterns.

For (8.21), the polar target is covered. At a valuation satisfying
the basis and B, no second-block answer can be true: such an answer
would make \(D^+_i\) true and the second conditional give \(\neg B\).
The coverage premise \(D_n\) therefore forces a first-block answer.
At a valuation satisfying the basis and \(\neg B\), the first
conditional similarly excludes every first-block answer, and coverage
forces one in the second block. These are the two required proper
witnesses. The first premise is used as a conditional, not as a
standalone asserted disjunction. When \(n=2\), both blocks are
singletons, proving the source's regularity assertion in all four cases.

To verify the limits for \(n>2\), take the original answers as
independent atoms and B,C as fresh atoms where they occur. One of the
two blocks has at least two members. For each chosen index j in the
first block, make only \(A_j\) true and set \(B=1,C=0\); for each
index j in the second block, make only \(A_j\) true and set
\(B=0,C=1\). Omit C in the two polar patterns. Every such
valuation satisfies each pattern's actual basis: (8.18)'s conditional
selects the matching block; (8.19)'s two conditionals and \(B\lor C\)
hold; (8.20)'s biconditionals match the block values; (8.21)'s two
conditionals and \(D_n\) hold. The response selecting the larger
block is fixed in all its valuations. For any proposed singleton
original answer in that block, choose a different index within it;
for one outside it, any of its valuations suffices. That response
therefore entails no individual original answer. This proves failure
of regularity in general, without a claim that every dependent
formula instance with \(n>2\) fails regularity.
\end{proof}
\end{annexitem}

\begin{annexitem}{annex:p18.ep.arbitrary-formulas}{thm:p18-ep-arbitrary-formulas}{Closed substitution for arbitrary formulas and all four answers}
\begin{proof}
The required closed-substitution property is established for arbitrary
formulas, not only predicates. On an atomic predicate or equality,
substituting c for a free x evaluates that occurrence at c's denotation.
The Boolean induction steps preserve this identity by their truth clauses.
Under a quantifier binding x, those bound occurrences are not replaced,
and evaluation uses the quantified individual's assignment to x on both
sides. Under a quantifier binding a different variable y, evaluate the
subformula at every y-value and apply its induction hypothesis, keeping
x at c's denotation. Since c has no free variables, no capture is
introduced. The universal and existential clauses therefore preserve
the property. This structural induction proves the closed-instance
identity for Ax,Bx,Cx with any permitted bound-variable occurrences.

Fix any model satisfying the universal premise. Evaluate it at c's
denotation and apply that identity to obtain equivalence between
\(A(x/c)\) and \(B(x/c)\land C(x/c)\). At every model exactly
one of the four binary conjunctions of B- and C-instances is true,
so the target is covered with either initial polar answer. For the
positive-positive target answer, the premise entails \(A(x/c)\).
For the positive-negative, negative-positive, and negative-negative
target answers, it entails \(\neg A(x/c)\). Thus each target
answer has the corresponding singleton original witness, nonempty
and proper in the two-answer original set.

Only evaluation at this particular named individual was used.
No step asserts that every domain individual is named, that the
three formulas are atomic, or that their other bound variables range
over a different model. These complete the arbitrary-formula extent.
\end{proof}
\end{annexitem}

\begin{annexitem}{annex:p18.ep.wh-polar}{thm:p18-ep-wh-polar}{Every wh-instance witness and precise lost-premise countermodels}
\begin{proof}
Polar target coverage holds at every model, hence with any original
wh-answer and the full basis. Its positive answer is itself an
original named-instance answer, giving a singleton witness.
Its negative answer \(\neg A(x/c_i)\), together with the existential
premise, supplies an A-individual. The restriction makes that individual
equal a listed c-denotation. It cannot equal the i-th denotation:
the substitution identity just proved would then make \(A(x/c_i)\)
true. Thus it equals some \(c_j\) with \(j\ne i\), and the same
identity gives a true answer \(A(x/c_j)\). This proves the stated
multiple-answer witness, without assuming distinct denotations.

Because x actually occurs free in Ax, substituting distinct constant
symbols produces syntactically distinct instances: at any one of those
free positions the substituted symbol differs. Hence the negative
response's witness set has \(n-1\) elements, is nonempty, and excludes
the original answer \(A(x/c_i)\). It is a proper subset of the whole
wh-answer set even if that set contains infinitely many other names.
For \(n=2\) it is a singleton, proving regularity.

For failure in general when \(n>2\), let Ax be the atomic Px.
Fix any proposed individual original answer Pc. Interpret distinctly
the finite set of needed symbols consisting of the list and c.
Choose a listed j different from i and from c if c is listed; this
is possible since at least two listed choices differ from i. Give
P extension just the j-th value. The two basis formulas and target
\(\neg Pc_i\) are true, but Pc is false. Other constants may be
assigned arbitrary existing values. This construction applies to
every candidate original answer, refuting the regularity clause.

If the existential premise is removed, take any nonempty named domain
with empty P-extension. The restriction and target \(\neg Pc_i\)
are true but every original answer is false. Thus no proper witness
can work. If the restriction is removed, take a domain with a further
element outside all constant denotations, interpret every constant
at named elements, and give P extension just that unnamed element.
The existential and negative target are true but all wh-answers are
false. The same second-clause failure occurs. These countermodels
show why the full basis, rather than an unrestricted bare existential
or a negative response by itself, establishes the reduction.
\end{proof}
\end{annexitem}
